% file thesis.tex
% Archivo thesis.tex
% Documento maestro que incluye todos los paquetes necesarios para el documento
% principal.

% Documento obtenido por un sinfin de iteraciones de administradores del LDC
% Estructura actual hecha por:
% Jairo Lopez <jairo@ldc.usb.ve>
% Actualizado ligeramente por:
% Alexander Tough 

\documentclass[oneside,11pt,letterpaper]{report}
\tolerance=1000  
\hbadness=10000  
\raggedbottom

% Paquetes para manejar graficos
\usepackage{epsf}
\usepackage[pdftex]{graphicx}
\usepackage{epsfig}
% Simbolos matematicos
\usepackage{latexsym,amssymb}
% Paquetes para presentar una tesis decente.
\usepackage{setspace,cite} % Doble espacio para texto, espacio singular para
                           % los caption y pie de pagina

% Paquetes no utilizados para citas
%\usepackage{mcite} 
%\usepackage{draft} 

\usepackage{wrapfig}
\usepackage{alltt}

% Acentos 
\usepackage[spanish,activeacute]{babel}
\usepackage[spanish]{translator}
\usepackage[utf8]{inputenc}
\usepackage{color, xcolor, colortbl}
\usepackage{multirow}
\usepackage{subfig}
\usepackage[OT1]{fontenc}
\usepackage{tocbibind}
\usepackage{anysize}
\usepackage{listings} 

% Para poder tener texto asiatico
%\usepackage{CJK}

% Opciones para los glosarios
\usepackage[style=altlist,toc,numberline,acronym]{glossaries}
\usepackage{url}
\usepackage{amsthm}
\usepackage{amsmath}
\usepackage{fancyhdr} % Necesario para los encabezados
\usepackage{fancyvrb}
\usepackage{makeidx} % En caso de necesitar indices.
\usepackage{tikz}
\usepackage{float}
\usepackage{caption}
\DeclareCaptionFont{white}{\color{white}}
\DeclareCaptionFormat{listing}{\colorbox{white}{\parbox{\textwidth}{#1#2#3}}}
\lstset{
 		 language=Pascal,
         numbers=left,               % Ort der Zeilennummern
         numberstyle=\tiny,          % Stil der Zeilennummern
         numbersep=5pt,              % Abstand der Nummern zum Text
         tabsize=5,                  % Groesse von Tabs
         extendedchars=true,         %
         breaklines=true,            % Zeilen werden Umgebrochen
         keywordstyle=\bfseries,
         frame=b,
         showspaces=false,           % Leerzeichen anzeigen ?
         showtabs=false,             % Tabs anzeigen ?
         xleftmargin=17pt,
         framexleftmargin=17pt,
         framexrightmargin=5pt,
         framexbottommargin=4pt,
         backgroundcolor=\color{white},
         showstringspaces=false      % Leerzeichen in Strings anzeigen ?        
}
\floatstyle{boxed} 
\restylefloat{figure}

\makeindex  % Necesitado para los indices

% Definiciones para definicions, teoremas y lemas
\theoremstyle{definition} \newtheorem{definicion}{Definici\'{o}n}
\theoremstyle{plain} \newtheorem{teorema}{Teorema}
\theoremstyle{plain} \newtheorem{lema}{Lema}

% Para la creacion de los pdfs
\usepackage{hyperref}

% Para resolver el lio del Unicode para la informacion de los PDFs
% En pdftitle coloca el nombre de su proyecto de grado/pasantia.
% En pdfauthor coloca su nombre.
\hypersetup{
    pdftitle = {ANÁLISIS DE METAHEURÍSTICAS PARA RESOLVER EL PROBLEMA DE ENRUTAMIENTO DE VEHÍCULOS CON ENVÍO Y RECEPCIÓN SIMULTÁNEA (VRPSPD)},
    pdfauthor={ALESSANDRO CHACÓN, RICARDO SANTANA},
    colorlinks,
    citecolor=black,
    filecolor=black,
    linkcolor=black,
    urlcolor=black,
    backref,
    pdftex
}

% Crea el glosario
\makeglossaries

% Incluye el glosario
\newacronym{as-h}{as-H}{proceso autosimilar con par\'{a}metro autosimilar $H$}

\newacronym{asie-h}{asie-H}{proceso autosimilar con par\'{a}metro autosimilar
$H$ e incrementos estacionarios}

\newacronym{aas-h}{aas-H}{proceso asint\'{o}ticamente autosimilar con
par\'{a}metro autosimilar $H$}


% Para crear la hoja escaneada de las firmas
\usepackage[absolute]{textpos}

% Pone los nombres y las opciones para mostrar los codigos fuentes
\newcommand{\vrp}{Problema de Enrutamiento de Vehículos }
\newcommand{\vrpingles}{Vehicle Routing Problem }
\newcommand{\vrpspd}{Problema de Enrutamiento de Vehículos con Entrega y Recepción Simultánea }
\renewcommand\lstlistingname{Algoritmo}
\renewcommand\lstlistlistingname{Lista de Algoritmos}

% Dimensiones de la pagina
\setlength{\headheight}{15pt}
\marginsize{2cm}{2cm}{2cm}{2cm}

% Se pueden omitir para que no compile ciertos capitulos.
\includeonly{header, resultados, conclusiones}

%%%%%%%%%%%%%%%%%%%%%%%%%%%%%%%%%%%%%%%%%%%%%%%%%%%%%%%%%%%%%%%%%%%%%%%%%%%
%%%%%%%%%%%%%%%%      end of preamble and start of document     %%%%%%%%%%%
%%%%%%%%%%%%%%%%%%%%%%%%%%%%%%%%%%%%%%%%%%%%%%%%%%%%%%%%%%%%%%%%%%%%%%%%%%%
\begin{document}

% Pagina de titulo
% Pagina de titulo
\begin{titlepage}
\begin{center}

% Upper part (aqui ya esta incluido el logo de la USB).
\includegraphics[scale=0.5,type=png,ext=.png,read=.png]{figures/cebolla} \\

% Encabezado
\textsc {\large UNIVERSIDAD SIMÓN BOLÍVAR} \\
\textsc{\bfseries DECANATO DE ESTUDIOS PROFESIONALES\\
COORDINACIÓN DE INGENIERÍA DE LA COMPUTACIÓN}

\bigskip
\bigskip
\bigskip
\bigskip
\bigskip
\bigskip
\bigskip
\bigskip
\bigskip

% Title/Titulo
% Aqui ponga el nombre de su proyecto de grado/pasantia larga
\textsc{\bfseries EXTRACCIÓN DE INFORMACIÓN FOCALIZADA BASADA EN RESPUESTAS INCOMPLETAS (FIE)}
\bigskip
\bigskip
\bigskip
\bigskip
\bigskip

% Author and supervisor/Autor y tutor
\begin{minipage}{\textwidth}
\centering
Por: \\
FRANCISCO RODRÍGUEZ DRUMOND \\

\bigskip
\bigskip
\bigskip

Realizado con la asesoría de: \\
PROF. SORAYA ABAD
\end{minipage}

\bigskip
\bigskip
\bigskip
\bigskip
\bigskip
\bigskip
\bigskip
\bigskip
\bigskip

% Bottom half
{PROYECTO DE GRADO \\ Presentado ante la Ilustre Universidad Simón Bolívar \\
como requisito parcial para optar al título de \\ Ingeniero de Computación} \\

\bigskip
\bigskip
\vfill

% Date/Fecha 
{\large \bfseries Sartenejas, Septiembre de 2012}

\end{center}
\end{titlepage}


% Pagina de acta final (vacio)
%% Pagina del acta final
\begin{titlepage}
\begin{center}

% Upper part
\includegraphics[scale=0.5,type=png,ext=.png,read=.png]{figures/cebolla} \\

\textsc {\large UNIVERSIDAD SIM'ON BOL'IVAR} \\
\textsc{DECANATO DE ESTUDIOS PROFESIONALES\\
COORDINACI'ON DE INGENIER'IA DE LA COMPUTACI'ON}

\bigskip
\bigskip
\bigskip
\bigskip
\bigskip
\bigskip

% Title
\textsc{ACTA FINAL PROYECTO DE GRADO}

\bigskip
\bigskip

% Aqui coloca el nombre de su proyecto de grado/pasantia larga.
\textsc{\bfseries  EXTRACCIÓN DE INFORMACIÓN FOCALIZADA BASADA EN RESPUESTAS INCOMPLETAS (FIE)}

\bigskip
\bigskip
\bigskip
\bigskip

\begin{minipage}{\textwidth}
\centering
Presentado por: \\
% Aqui coloca su nombre.
\textsc{\bfseries FRANCISCO RODRÍGUEZ DRUMOND} \\

\bigskip
\bigskip
\bigskip
\bigskip

Este Proyecto de Grado ha sido aprobado por el siguiente jurado examinador: \\

\bigskip
\bigskip

% Despues de cada line coloca el (los) nombre(s) de
% cada uno de los integrantes del jurado.
\line(1,0){200} \\
PROF. 1\\

\bigskip
\bigskip

\line(1,0){200} \\
PROF. 2\\

\bigskip
\bigskip

\line(1,0){200} \\
PROF. 3\\
\end{minipage}

\bigskip
\bigskip
\vfill

% Date/Fecha
{\large \bfseries Sartenejas, X/X/12}

\end{center}
\end{titlepage}

\setcounter{secnumdepth}{3}
\setcounter{tocdepth}{4}

% Define encabezado numeros romanos y como se separan los captiulos y las
% secciones
\addtolength{\headheight}{3pt}
\pagenumbering{roman}
\pagestyle{fancyplain}

\renewcommand{\chaptermark}[1]{\markboth{\chaptername\ \thechapter:\,\ #1}{}}
\renewcommand{\sectionmark}[1]{\markright{\thesection\,\ #1}}

\onehalfspacing

\lhead{}
\chead{}
\rhead{}
\renewcommand{\headrulewidth}{0.0pt}
\lfoot{}
\cfoot{\fancyplain{}{\thepage}}
\rfoot{}


% Pagina de resumen
\setcounter{page}{4}
\chapter*{Resumen}

DIA es una arquitectura para realizar consultas sobre información presente en un conjunto de documentos. Para ello, contempla poblar una base de datos con información extraída en una búsqueda amplia sobre un conjunto de documentos de interés; posteriormente se pueden realizar preguntas sobre esa información utilizando una ontología. La extracción focalizada es una fase que contempla DIA para buscar información en documentos una vez que se haya determinado que falta información para poder responder alguna consulta. El objetivo de este trabajo fue diseñar e implementar una propuesta para la tarea de Extracción Focalizada. Para ello se estudió el problema de la determinación de los orígenes de incompletitud en consultas hechas en DIA, se razonó sobre la información presente en la ontología y en la base de datos que se pueda utilizar para encontrar la información faltante y se construyó un extractor de información que trabaja con expresiones regulares. La solución propuesta fue probada utilizando actas de Consejos Directivos y Académicos de la USB, para realizar extracción sobre designaciones de jurados de ascenso, ascenso en el escalafón y designaciones administrativas. Los resultados obtenidos fueron: \\

Máximo 250 palabras.

\newpage



% Pagina de dedicatoria (opcional)
\setcounter{page}{5}

\vspace*{8cm} 
\pdfbookmark[0]{Dedication}{dedication} % Sets a PDF bookmark for the dedication
\begin{center} 
\large DEDICATORIA
\end{center}
\newpage


% Pagina de agradecimientos (opcional)
\setcounter{page}{6}

\chapter*{Agradecimientos
\markboth{Agradecimientos}{Agradecimientos}}
AQU\'I VAN LOS AGRADECIMIENTOS.





% Crea la tabla de contenidos
\tableofcontents

% Crea la lista de cuadros
\listoftables

% Crea la lista de figuras
\listoffigures

% Crea la lista de codigos fuentes
%\lstlistoflistings

\clearpage

% Define encabezado en numeros arabicos  
\pagenumbering{arabic}

\fancyhf{} % Redefine el encabezado 
\lhead{}
\chead{}
\rhead{\fancyplain{}{\thepage}}
\renewcommand{\headrulewidth}{0.0pt}
\lfoot{}
\cfoot{}
\rfoot{}

%\doublespacing


% Incluye los archivos deseados - El contenido de
% su proyecto de grado/pasantia larga.
% Marco Teorico.
\chapter{El Problema de Enrutamiento de Vehículos} \label{chap:vrp}


La distribución de bienes es una tarea fundamental hoy en día para muchas empresas y organizaciones. Pro\-ble\-mas como conseguir la distancia más corta entre dos establecimientos, satisfacer demandas de clientes y mejorar el tiempo de la entrega son factores que influyen en el costo final de un servicio. Teniendo en cuenta lo importante que es la distribución para muchos aspectos en la sociedad los investigadores han estudiado formas de reducir los costos, haciendo la distribución más eficiente. El Problema de Enrutamiento de Vehículos (\textbf{VRP} por sus siglas en inglés de Vehicle Routing Problem) encapsula las restricciones y definiciones de la distribución de bienes en su forma más general. El problema fue propuesto originalmente por Dantzig y Ramser en 1959 \cite{primervrp} mientras intentaban encontrar maneras más eficientes de entregar gasolina a estaciones de servicio utilizando camiones como medio de transporte. \\

VRP está formado de los siguientes elementos:
\begin{itemize}
	
	\item Un depósito que almacena los bienes.
	\item Un conjunto de clientes que demandan una cantidad exacta de bienes.
	\item Un conjunto de vehículos cuyo objetivo es transportar los bienes entre el depósito y los clientes.

\end{itemize}

La resolución exacta del problema es encontrar las rutas que deben recorrer los vehículos para satisfacer las demandas de los clientes de tal manera de minimizar el costo (o distancia) de las rutas.

\section{Descripción} \label{sect:descripcion}

Se tiene un depósito que contiene una cantidad finita de bienes, los cuales se deben distribuir utilizando una cantidad de vehículos. Éstos deben recorrer los clientes satisfaciendo sus demandas de productos cumpliendo las siguientes restricciones:

\begin{itemize}

	\item Todos los vehículos deben comenzar su recorrido en el depósito y teminar en el depósito.
	\item Cada cliente debe ser visitado exactamente por un sólo vehículo una única vez.
	\item Se deben satisfacer las demandas de todos los clientes en su totalidad.
	
\end{itemize}

La resolución del problema consiste en minimizar el costo total de todos los recorridos de los vehículos cum\-plien\-do las restricciones mencionadas anteriormente. \\


VRP pertenece al conjunto de problemas \textit{NP-Hard
}\footnote{NP-hard: una clase de problemas de decisión que pueden ser resueltos por una máquina de turing no determinista en tiempo polinomial.}. Debido a esto, la utilización de métodos exactos resulta inconveniente para resolver el problema, ya que su tiempo de resolución incrementa de forma exponencial a medida de que las instancias aumentan de tamaño (número de clientes). Generalmente para resolver este tipo de problemas se utilizan algoritmos que logran resultados aproximados de buena calidad con tiempos reducidos. Estos algoritmos son llamados metaheurísticas y han demostrado ser de gran ayuda, siendo así un área de investigación de gran interés dentro de la computación.

Más adelante se describirán las metaheurísticas utilizadas.

\begin{comment}
\section{Formulación} \label{sect:formulacion}

VRP puede ser definido formalmente utilizando la formulación propuesta por Level y Palhazi\cite{tesisdanielernesto}: \\

Usualmente el problema es modelado utilizando un grafo, donde los vértices representan los clientes (incluyendo el depósito) y los arcos llevan asociados un costo (que usualmente es visto como la distancia entre dos puntos).

Sean $n$ el número de clientes, $k$ el número de vehículos disponibles y $G = (V, E)$ un grafo completo donde $V = \{0,1,...,n\}$ es el conjunto de vértices y $E = \{(i,j)|(i,j) \in V \times V \ \wedge \ i \not= j\}$ es el conjunto de arcos. El vértice $0$ representa al depósito y los vértices $q = 1,2,...,n$ corresponden a los clientes que deben ser suplidos. Los arcos representan los costos de los caminos que puede utilizar la red de transporte.

Una solución esta formada por una conjunto de rutas $R = \{r_1,r_2,...,r_k\}$, las cuales deben cumplir las siguientes condiciones:

\begin{itemize}

	\item Cada ruta $r_i \in R$ debe iniciar y terminar en el depósito.
	\item Cada vértice $v_i \in V - \{0\}$ debe ser visitado una y sólo una vez.

\end{itemize}
\end{comment}


\section{Variantes de VRP} \label{sect:variantes}

En la vida real muchos problemas de distribución han tratado de modelarse utilizando VRP, sin embargo resulta complejo, por lo que ha surgido la necesidad de crear nuevas formas de modelado más específicas de acuerdo a cada tipo de problema, en algunos casos agregando restricciones y en otros relajando algunas de ellas. Estas nuevas formas de modelado son llamadas variantes. A continuación se mencionan algunas:

El Problema de Enrutamiento de Vehículos con Capacidad (CVRP por sus siglas en inglés de Capacited Vehicle Routing Problem) se define por contar con una flota de vehículos homogéneos que poseen una capacidad máxima fija. Este elemento adicional crea una nueva restricción al problema: no permitir que los bienes dentro del vehículo excedan su capacidad máxima en ningún momento.

En el \vrp con Ventanas de Tiempo (VRPTW por sus siglas en inglés de \vrpingles with Time Windows) se define para cada cliente un rango de tiempo o \emph{ventana de tiempo}, durante el cual debe ser atendido por algún vehículo. Por lo tanto un vehículo puede atender a un cliente si y sólo si el vehículo llega al cliente dentro su ventana de tiempo.

El \vrp con Recolección (VRPB por sus siglas en inglés de Vehicle Routing Problem with Backhauls) requiere que los clientes se clasifiquen en dos conjuntos disjuntos: clientes que reciben bienes del depósito (llamados consumidores o \textit{linehauls}) y clientes que entregan bienes al depósito (llamados pro\-ve\-e\-do\-res o \textit{backhauls}). En esta variante se requiere que los consumidores sean atendidos antes que los proveedores, permitiendo así asegurar que los vehículos nunca superen su capacidad máxima, entregando todo lo que necesitan entregar antes de recoger nuevos bienes.

El \vrp Estocástico (SVRP por sus siglas en inglés de Stochastic Vehicle Routing Problem) se caracteriza por poseer uno o varios componentes aleatorios. Podrían haber nuevos clientes durante la ejecución, las demandas podrían ser aleatorias y conocidas sólo al momento de que un vehículo llegue al cliente.

%El \vrp Dinámico (DVRP por las siglas en inglés de Dynamic Vehicle Routing Problem) es una variante compleja ya que  mientras se construye la ruta aún no se sabe información completa de las demandas de ciertos clientes. Además, luego de construidas las rutas la información de las demandas podría cambiar dinámicamente, por lo tanto mientras los vehículos se desplazan podría haber cambios que influyan positivamente o negativamente en el costo total de las rutas. 

El Problema de Enrutamiento de Vehículos con Múltiples Depósitos (VRPMD por sus siglas en inglés de \vrpingles with Multiple Depots) agrega depósitos adicionales al problema general. Por lo tanto una flota de vehículos distinta se encarga de atender a los clientes.

El Problema de Enrutamiento de Vehículos con Múltiples Vehículos (VRPMV por las siglas en inglés de Vehicle Routing Problem with Multiple Vehicles) tiene como elemento clave que la distribución de los bienes no esté limitada a un sólo vehículo por cliente, ya que en algunas situaciones la demanda de bienes de parte de un cliente es mayor a la capacidad máxima del vehículo. Elimina la restricción de que un cliente sea atendido por sólo un vehículo.

A continuación se pasa a describir una variante de particular interés de este proyecto, como lo es VRPSPD.

%\section{Aplicaciones de VRP} \label{sect:aplicaciones}
%
%En los sectores productivos es muy común observar la necesidad de distribuir bienes, así como recolectar materiales. Cada día aumenta más la cantidad de clientes en una cadena de distribución, por lo que ha sido necesario desarrollar formas más eficientes de realizar entregas y recolecciones óptimas con el fin de disminuir los costos. El \vrp se ha investigado con mucho hincapié para resolver este problema. La utilización de algoritmos computacionales ha resultado, a medida que se profundizan las investigaciones, beneficiosa para reducir los costes de distribución en las organizaciones.
%
%VRP ha sido modelado matemáticamente con el objetivo de que pueda ser representado computacionalmente. Esta representación ha llevado a muchas empresas e industrias a llevar a cabo implementaciones en máquina que resuelvan VRP y así contribuir a minimizar los costos de distribución. Algunas aplicaciones fueron implementadas por algunos investigadores:
%
%\begin{description}
%\item[Golden y Wasil\cite{aplicacionvrp1}: ] Utilizaron un software para optimizar la distribución de la industria de bebidas. Tanto de ventas preordenadas por los clientes, como entrega de bebidas compradas por los clientes al conductor del vehículo. Estudios concluyeron en que en Estados Unidos el 49.8\% de las ventas fueron realizadas directamente por el chofer, lo que implica que los costos asociados a la distribución de las bebidas eran enormes. Se realizaron implementaciones de tres casos de estudio, buscando reducir los costos en un 2-5\%.
%\item[Evans y Norback\cite{aplicacionvrp2}: ] Implementaron un programa para reducir los costos de distribución de alimentos en la empresa Kraft Inc. Utilizaron un sistema basado en heurísticas utilizando gráficos de computadora para ilustrar las largas rutas de su cadena de distribución. En pruebas, el personal de Kraft logró reducir los costos operativos en 10.7\% en problemas de hasta 223 paradas en clientes.
%
%\end{desocription}
\section{VRP con Recepción y Entrega Simultánea (VRPSPD)} \label{sect:vrpspd}

El \vrpspd (VRPSPD por sus siglas en inglés de Vehicle Routing Problem with Simultaneous Pickup and Delivery) fue abordado por primera vez por Min \cite{primervrpspd} en el año 1989. En esta variante los clientes pueden, de manera simultánea, recibir y entregar bienes. Por lo tanto, cada cliente tiene asociado una cantidad de bienes por recibir y por entregar. 

\subsection{Modelo Matemático} \label{subsect:vrpspdFor}

\subsubsection{Características del Problema}

\begin{itemize}
\item Las demandas de entrega y recepción de bienes son determinísticas.
\item Las demandas de recepción son atendidas exclusivamente desde el depósito.
\item No hay intercambio de bienes entre clientes.
\item Las demandas de entrega y recepción de bienes no pueden ser divididas.
\item Todo cliente es visitado una sola vez.
\item No hay restricción de distancia máxima recorrida para los vehículos.

\end{itemize}

\subsubsection{Notación}

\begin{tabular}{ll}
{$n:\,$} 		& {Número total de clientes.}  \\ 
{$Q:\,$} 		& {Capacidad de los vehículos.} \\ 
{$k:\,$}  		& {Número de vehículos.}\\
{$i,j:\,$}		& {Índices para el deposito y los clientes. El depósito se denota con el número 0.}\\
{$l:\,$}    	& {Índice para los vehículos.}\\
{$c_{ij}:\,$} 	& {Distancia (costo) entre el cliente $i$ y el cliente $j$.}\\
{$p_j:\,$}		& {Bienes a ser entregados por un cliente $j$, donde $j=1,2,…,n$.}	\\	
{$d_j:\,$}		& {Bienes a ser recogidos por un cliente $j$, donde $j=1,2,…,n$.} \\
{$x_{ijl}$} 	& {Variable de decisión que toma el valor de 1 si el vehículo $l$ visita al cliente $i$ antes } \\
{$ $}              & {que al cliente $j$, de lo contrario es 0.} \\
{$y_{ij}$} 		& {Bienes recogidos hasta el cliente $i$ y transportados hasta el cliente $j$.} \\
{$z_{ij}$}		& {Bienes por entregar desde el cliente $i$ y transportados hasta el cliente $j$.}\\

\end{tabular} 

\subsubsection{Formulación Matemática}

La función objetivo ~\eqref{eq:FOBJ} correspondiente a esta formulación matemática es la siguiente: 


\begin{equation}\label{eq:FOBJ}
 Min \sum_{l=1}^k\sum_{i=0}^n\sum_{j=0}^n  c_{ij} x_{ijl} 
\end{equation}

Sujeto a:

\begin{equation}\label{eq:FM1}
\sum_{i=0}^n \sum_{l=1}^k x_{ijl}=1 , \: j=1,2,…,n 
\end{equation}


	La restricción ~\eqref{eq:FM1} asegura que cada cliente es visitado exactamente una vez por un sólo vehículo.


\begin{equation}\label{eq:Ft2}
\sum_{i=0}^n x_{ijl} - \sum_{i=0}^n x_{jil} = 0, \: j=0,1,…,n, l=1,2,…,k  
\end{equation}	
	
	La restricción ~\eqref{eq:Ft2} asegura que el mismo vehículo entra y sale de cada cliente que visita.


\begin{equation}\label{eq:FM3}
\sum_{j=1}^n x_{0jl} \leq 1 , \: l=1,2,…,k   
\end{equation}

	La restricción ~\eqref{eq:FM3} representa que la cantidad de vehículos utilizados es menor o igual que  $k$. 



\begin{equation}\label{eq:FM4}
\sum_{i=0}^n y_{ji} - \sum_{i=0}^n y_{ij} = p_j, \: \forall j \neq 0  
\end{equation}

\begin{equation}\label{eq:FM5}
\sum_{i=0}^n z_{ji} - \sum_{i=0}^n z_{ij} = d_j, \: \forall j \neq 0  
\end{equation}


Las restricciones ~\eqref{eq:FM4} y ~\eqref{eq:FM5} son ecuaciones de flujo para la demanda de recepción y de entrega de mercancía.


\begin{equation}\label{eq:FM6}
(y_{ij}+ z_{ij})  \leq Q \sum_{l=1}^k x_{ijl}   , \: i,j=0,1,…,n  
\end{equation}

La restricción ~\eqref{eq:FM6} establece que los bienes  de recepción y de entrega solo pueden ser transportados por arcos pertenecientes a la solución y para ninguno de ellos se excede de la máxima capacidad permitida.

\begin{equation}\label{eq:FM7}
x_{ijl}\in \lbrace0,1\rbrace,  y_{ij} \geq 0 ,z_{ij} \geq 0,   i,j=0,1,…,n,   l=0,1,…,k     
\end{equation}		

		Finalmente  la restricción ~\eqref{eq:FM7} muestra la naturaleza de las variables de decisión.

%	Existen  $k$ vehículos en el depósito, V representa al conjunto de clientes que serán visitados, donde $n = |V|$ es el número total de clientes. La ubicación del depósito y los clientes es conocida.  Todo cliente posee una cantidad conocida de mercancía demandada y de mercancía a devolver, representadas por $d_j$   y $p_j$ respectivamente con valores desde $j=1,2,…,n$.  Toda ruta debe comenzar en el depósito y debe terminar en el mismo. Todos los clientes son visitados exactamente una vez y sus necesidades son satisfechas por un solo vehículo. $V_0 = V \cup {0}$  , es el conjunto de clientes más el depósito; $c_{ij}$ es la distancia entre el cliente $i$ y el cliente $j$, y la capacidad del vehículo es representada por $Q$. $x_{ijl}$ representa la variable de decisión que toma el valor de 1 si el vehículo $l$ visita visita al cliente $i$ antes que al cliente $j$, de lo contrario es 0. $y_{ij}$  representa los bienes recogidos hasta el cliente $i$ y transportados hasta el cliente $j$. $z_{ij}$ Representa los bienes por entregar desde el cliente $i$ y transportados hasta el cliente $j$.


%Puedes quitar esto(es opcional)
\vspace{5 mm}



\chapter{Metaheurísticas} \label{chap:metaheuristicas}

\section{Definición} \label{sect:definicion}

Ya que el problema de enrutamiento de vehículos es un problema complejo, resolverlo exactamente implica un incremento exponencial del tiempo de cómputo a medida que las instancias aumentan de tamaño. Para problemas complejos usualmente se trata de conseguir una buena solución que tome un tiempo aceptable. Una forma de lograr esto es utilizando las denominadas \emph{metaheurísticas}. El objetivo de esta clase de algoritmos  es lograr un balance entre obtener buenos tiempos de ejecución y buenas soluciones a través de un proceso iterativo. La palabra \emph{heurística} proviene del griego $\varepsilon\upsilon\rho\iota\sigma\kappa\varepsilon\iota\nu$, que significa \emph{``hallar, encontrar''}. El prefijo \emph{meta} quiere decir \emph{``nivel superior''}. Juntas podrían interpretarse como `mejorar lo encontrado'.

Las metaheurísticas son una familia de algoritmos aproximados de propósitos generales y no determinísticos que consiste de procesos iterativos que guían una heurística subordinada, realizando sobre el espacio de búsqueda exploración y explotación.
\begin{itemize}
\item \textbf{Exploración:} Busca en el espacio de soluciones lugares nunca antes visitados. 
\item \textbf{Explotación:} Regresa a lugares del espacio de soluciones ya visitados para obtener soluciones anteriormente calificadas como buenas.
\end{itemize}

\section{Clasificación} \label{sect:clasificacion}
Una metaheurística puede ser:

\begin{itemize}
\item \textbf{Constructiva:} Construye paso a paso una solución del problema, basándose en la mejor elección de cada iteración.
\item \textbf{De trayectoria:} Parten de una solución inicial y la mejoran progresivamente hasta lograr el óptimo.
\item \textbf{Poblacional:} Basado en conjunto de soluciones que evolucionan sobre el espacio de búsqueda.
\end{itemize}

\section{Metaheurísticas de interés} \label{sect:interes}

A continuación una breve descripción de las metaheurísticas más relevantes para este trabajo.

\subsection{Búsqueda Local} \label{subsect:ls}

La \emph{búsqueda local} (LS por sus siglas en inglés de \emph{Local Search}) es uno de los métodos más empleados para resolver problemas complejos. Su orígen proviene del año 1950 donde fue utilizada por primera vez para resolver el problema TSP.\footnote{Siglas en inglés de Traveling Salesman Problem, problema en donde un vendedor debe visitar todos los clientes una única vez, reduciendo su desplazamiento al mínimo} 
\\
Como descrito en \cite{LocalSearchVNDDef} el algoritmo parte de una solución inicial, que se representa en la solución actual, de la cual se obtiene una vecindad parcial o completa de nuevas soluciones. A partir de esta vecindad se elige una nueva solución actual, siempre y cuando  la solución sea mejor que la solución actual de acuerdo a un criterio previamente establecido. Este proceso se repite hasta alcanzar un criterio de parada previamente definido. 

El pseudocódigo general de LS se puede ver en el \textbf{Algoritmo~\ref{alg:LSS}}.

\begin{code}[includerangemarker=false,frame=single,label=alg:LSS,caption=Pseudocódigo de Búsqueda Local,firstnumber=100, mathescape]
$S_0$ := ConstruirSolucionInicial()
$S_{actual}$ := $S_0$

while ($\neg$criterio_parada) do
	vecindad := ConstruirVecindad($S_{actual}$)
	Seleccionar $S$ $\in$ vecindad | $S$ := Min{Costo($S$)}
	if (Costo($S$) $<$ Costo($S_{actual}$))	
		$S_{actual}$ := S
	end if		
end while

$\textbf{return}$ $S_{actual}$
\end{code}

\subsection{Búsqueda Local Iterada} \label{subsect:ils}

La \emph{búsqueda local iterada} (ILS por sus siglas en inglés de \textit{Iterated Local Search}) es una metaheurística de trayectoria que consiste en iterativamente perturbar la solución actual a la cual se le aplica una heurística de mejoramiento. Una perturbación consiste en aplicar un operador de manera aleatoria, obteniendo de esta forma una nueva solución. La fuerza de una perturbación se refiere al número de componentes de la solución que son modificadas, es importante destacar que si la perturbación es muy fuerte el ILS puede comportarse como si su solución actual fuera aleatoria, mientras que si la perturbación es muy leve la solución actual normalmente caerá de nuevo en el mismo óptimo local, dejando un espacio de diversificación muy limitado \cite{ILSDef}. 
\\
El pseudocódigo general de ILS se puede ver en el \textbf{Algoritmo~\ref{alg:ILS}}.

\begin{code}[includerangemarker=false,frame=single,label=alg:ILS,caption=Pseudocódigo de Búsqueda Local Iterada,firstnumber=100, mathescape]
$S_0$ := ConstruirSolucionInicial()
$S_{mejor}$ := LocalSearch($S_0$)
while ($\neg$criterio_parada) do

	$S_{perturb}$ := Perturbar($S_{mejor}$)
	$S_{actual}$ := LocalSearch($S_{perturb}$)

	if(Costo($S_{mejor}$) $>$ Costo($S_{actual}$))		
		$S_{mejor}$ := $S_{actual}$		
	end if
			
end while

$\textbf{return}$ $S_{mejor}$
\end{code}

\subsection{Búsqueda en Vecindades Variables} \label{subsect:vnd}

La \emph{búsqueda en vecindades variables} (VNS por sus siglas \emph{Variable Neightborhood Search}) es una metaheurística de trayectoria que consiste en cambiar sistemáticamente de vecindad dentro de una búsqueda local, además de un método de perturbación para diversificar la solución. Una de sus variantes, VND (\emph{Variable Neightborhood Descent} o \emph{Búsqueda Descendente en Vecindades Variables}) se obtiene cuando los cambios de vecindarios son realizados de manera determinística. 
\\
El pseudocódigo general de VND se puede ver en el \textbf{Algoritmo~\ref{alg:VNS}} \cite{VNDDef}.

\begin{code}[includerangemarker=false,frame=single,label=alg:VNS,caption=Pseudocódigo de Búsqueda Descendente en Vecindades Variables,firstnumber=100, mathescape]
Seleccionar vecindades $N_{k}$ para 
k := 1,...,$K_{max}$ 
$S_0$ := ConstruirSolucionInicial()
$S_{mejor}$ := $S_0$
while ($\neg$criterio_parada) do
    k := 1
    while (k < $k_{max}$+1) do 
    	$vecindad$ := ConstruirVecindadVND($N_{k}$,($S_{mejor}$))
		Seleccionar $S_{actual}$ $\in$ $vecindad$ | $S_{actual}$ := Min{Costo($S_{actual}$)}	
		if(Costo($S_{mejor}$) > Costo($S_{actual}$))		
			$S_{mejor}$ := $S_{actual}$
			k := 1
		else
		 k := k+1
		end if 			
	end while
end while

$\textbf{return}$ $S_{mejor}$
\end{code}

\subsection{Búsqueda Tabú} \label{subsect:ts}

La \emph{búsqueda tabú} (TS por sus siglas en inglés de \emph{Tabu Search}) es una metaheurística de trayectoria propuesta por Fred Glover en 1986 \cite{tabusearch} con el objetivo de ayudar a la \emph{búsqueda local} a evitar el estancamiento en óptimos locales, permitiendo movimientos dentro del espacio de soluciones que no necesariamente impliquen un mejor valor de función objetivo que el valor actual.
Para evitar el fenómeno cíclico de la búsqueda local y así impedir volver a visitar soluciones anteriormente visitadas, éstas se prohiben mediante la utilización de memoria a corto plazo, empleando la llamada \emph{lista tabú}, cuyo propósito es almacenar la historia reciente de la búsqueda. 
Estas soluciones ya visitadas son declaradas \emph{tabú}, es decir, su posterior exploración está prohibida por un determinado número de iteraciones definidas por el valor de la \emph{tenencia tabú}. 

Usualmente es impráctico guardar soluciones en la lista tabú por razones de uso excesivo de recursos. Por este motivo normalmente se guardan movimientos específicos entre soluciones permitiendo así un uso eficiente de recursos y buenos resultados.

Suele ocurrir que la búsqueda tabú evite movimientos que dirijan a muy buenas soluciones como resultado de la utilización de la lista tabú, por lo que generalmente se define un \emph{criterio de aspiración}, que permite a la búsqueda ignorar el hecho de que una solución se encuentre en la lista tabú sólo si esa solución conlleva un valor de función objetivo mejor que el de la mejor solución encontrada hasta el momento.

El pseudocódigo general de TS se puede ver en el \textbf{Algoritmo~\ref{alg:TS}}:

\begin{code}[includerangemarker=false,frame=single,label=alg:TS,caption=Pseudocódigo de Búsqueda Tabú,firstnumber=100, mathescape]
$S_0$ := ConstruirSolucionInicial()
$S_{actual}$ := $S_0$
$S_{mejor}$ := $S_0$
$Lista\_Tabu$ := $\emptyset$

while ($\neg$criterio_parada) do
	$vecindad$ := ConstruirVecindad($S_{actual}$)
	Seleccionar $S$ $\in$ $vecindad$ | $S$ = Min{Costo($S$)}
	if (Costo($S$) < Costo($S_{mejor}$))	        // criterio de aspiracion
		$S_{mejor}$ := $S$
	else
		if ($S$ $\in$ $Lista\_Tabu$)
			Seleccionar $S$ $\in$ $vecindad$ | $S$ := Min{Costo($S$)} $\wedge$ $S$ $\notin$ $Lista\_Tabu$
		end if
	end if
	$S_{actual}$ := $S$
	Guardar $S_{actual}$ en $Lista\_Tabu$ y actualizar tenencia si necesario
end while

$\textbf{return}$ $S_{mejor}$
\end{code}

\subsection{Búsqueda Local Guiada} \label{subsect:gls}

La \emph{búsqueda local guiada} (GLS por sus siglas en inglés de \emph{Guided Local Search}) es una metaheurística de trayectoria diseñada para guiar la \emph{búsqueda local} penalizando características no deseables de la solución actual, para así evitar el estacamiento en óptimos locales. Primeramente se definen las características de una solución candidata, cuando la búsqueda se atasca en un óptimo local las soluciones que posean  ciertas características consideradas no deseables son penalizadas. Usualmente, se utiliza una función de utilidad tal que la característica que maximice la función será penalizada. La característica penalizada y el número de veces que ha sido penalizada será almacenada de manera que en búsquedas posteriores el valor de la función objetivo se vea alterado. De esta manera, la búsqueda se concentra en explorar mayormente soluciones calificadas como deseables y permite prevenir que dirija todo su esfuerzo en explorar sólo una región del espacio de soluciones.

El pseudocódigo general de GLS se puede ver en el \textbf{Algoritmo~\ref{alg:GLS}}:

\begin{code}[includerangemarker=false,frame=single,label=alg:GLS,caption=Pseudocódigo de Búsqueda Local Guiada,firstnumber=100, mathescape]
Funcion $U$
Conjunto $C$

$S_0$ := ConstruirSolucionInicial()
$S_{actual}$ := $S_0$
$S_{mejor}$ := $S_0$

while ($\neg$criterio_parada) do
	$S_{mejor}$ := LocalSearch($S_{actual}$)
	Encontrar la caracteristica $C_i$ | $C_i$ = Max {$U(i)$}
	Penalizar($C_i$)
	$S_{actual}$ := $S_{mejor}$
end while

$\textbf{return}$ $S_{mejor}$
\end{code}

\subsection{Optimización Basada en Colonia de Hormigas} \label{subsect:aco}

La \emph{optimización basada en colonia de hormigas} (ACO por sus siglas en inglés de \emph{Ant Colony Optimization}) es una metaheurística poblacional inspirada en los rastros de feromonas dejados por las hormigas al caminar y el comportamiento de ellas ante las feromonas. Estas feromonas son utilizadas por las hormigas como medio de comunicación con el fin de conseguir una forma óptima para llegar a la comida. Análogamente ACO está basado en la comunicación indirecta de las hormigas artificiales mediante los rastros de feromonas artificiales. 

Se tiene un número fijo de hormigas, cada hormiga construye una solución propia basándose en información subyacente a caminos de feromonas dejados por otras hormigas. Si no existen feromonas en los caminos a tomar, se escoje un camino aleatorio. Si existen feromonas en los caminos a tomar, se elige probabilísticamente dependiendo de la cantidad de feromonas depositada en los caminos. Mientras más feromonas tenga un camino mayor probabilidad tiene de ser elegido. Cada hormiga deposita feromonas en su recorrido. 

Para evitar que el flujo de hormigas se dirija siempre por el camino de más feromonas, acarreando un es\-tan\-ca\-mien\-to en un óptimo local, se define un factor de evaporación, que reduce la cantidad de  feromonas en los caminos cada cierta cantidad de iteraciones o cada cierto tiempo. De esta manera se explorará eficazmente el espacio de soluciones, logrando soluciones de calidad.

El pseudocódigo general de ACO se puede ver en el \textbf{Algoritmo~\ref{alg:ACO}}:

\begin{code}[includerangemarker=false,frame=single,label=alg:ACO,caption=Pseudocódigo de Optimización basada en Colonia de Hormigas,firstnumber=100, mathescape]
hormiguero := CrearHormigas()
InicializarFeromonas()
$S_{mejor}$ := ConstruirSolucionInicial()

while ($\neg$criterio_parada) do
	$\textbf{foreach}$ $hormiga \in$ hormiguero do
		$S_{actual}$ := ConstruirSolucion($hormiga$)
		ActualizarFeromonas($S_{actual}$)
		if (Costo($S_{actual}$) < Costo($S_{mejor}$))
			$S_{mejor}$ := $S_{actual}$
		end if
	end $\textbf{foreach}$	
	EvaporarFeromonas()
end while

$\textbf{return}$ $S_{mejor}$
\end{code}

\subsection{Optimización por Enjambre de Partículas} \label{subsect:pso}

La \emph{optimización por enjambre de partículas} (PSO por sus siglas en inglés de \emph{Particle Swarm Optimization}) es una metaheurística poblacional inspirada en el comportamiento en grupo de enjambres. PSO imita los movimientos físicos individuales en los enjambres como método de búsqueda. 

Un enjambre está compuesto por un número definido de partículas, las partículas tienen dos características esenciales: posición y  velocidad. La posición de una partícula representa una solución al problema, mientras que la velocidad representa la habilidad de la partícula  para cambiar de posición. Utilizando estos dos factores, se intenta encontrar la mejor posición posible dentro de un espacio de búsqueda, la  cual se espera represente una buena solución al problema.

La velocidad de una partícula es constantemente actualizada basándose en tres términos: (1) Inercia, obliga a la partícula a moverse en la misma dirección. (2) Aprendizaje cognitivo, obliga a la partícula a moverse a la mejor posición encontrada hasta el momento por ella. (3) Aprendizaje social, obliga a la partícula a moverse a la mejor posición encontrada hasta el momento por todas las partículas del enjambre.
La velocidad de cada partícula es influenciada por los factores cognitivos y sociales con el fin de dirigir el enjambre a encontrar mejores soluciones al problema. Además, se incluye un factor estocástico para evitar que todas las partículas con misma posición se muevan en las mismas direcciones y así explorar extensivamente el espacio de soluciones.

El pseudocódigo general de PSO se puede ver en el \textbf{Algoritmo~\ref{alg:PSO}}:

\begin{code}[includerangemarker=false,frame=single,label=alg:PSO,caption=Pseudocódigo de Optimización por Enjambre de Partículas,firstnumber=100, mathescape]
$\emph{enjambre}$ := CrearEnjambre()
InicializarPosiciones($\emph{enjambre}$)
InicializarVelocidades($\emph{enjambre}$)
while ($\neg$criterio_parada) do
	$\textbf{foreach}$ $\emph{p}$ $\in$ $\emph{enjambre}$ do
		ActualizarVelocidad($\emph{p}$)
		ActualizarPosicion($\emph{p}$)
		if (Costo($\emph{p}$) < MejorCosto($\emph{p}$))
			ActualizarMejorCosto($\emph{p}$)
		end if
		if (Costo($\emph{p}$) < MejorCosto($\emph{enjambre}$))
			mejor_particula := $\emph{p}$
		end if
	end $\textbf{foreach}$
end while
$S_{mejor}$ := ObtenerPosicion(mejor_particula)

$\textbf{return}$ $S_{mejor}$
\end{code}

\subsection{Búsqueda en Dispersión} \label{subsect:sca}

La \emph{búsqueda en dispersión} (SS por sus siglas en inglés de \emph{Scatter Search}) es una metaheurística poblacional la cual está diseñada para operar a través de un conjunto de soluciones, llamado conjunto de referencia, el cual está constituido por buenas soluciones obtenidas hasta el momento. Es importante destacar que una buena so\-lu\-ción puede ser considerada de esta forma tanto por su diversidad, como por el valor de su función objetivo. Este enfoque genera sistemáticamente nuevas soluciones a partir del conjunto de referencia a través de combinaciones de soluciones pertenecientes al mismo.    
\\
El algoritmo básico SS consta de cinco métodos:
\\
\begin{enumerate}
 
\item Un \textit{Método de generación diversificada} que crea un conjunto de soluciones iniciales.

\item Un \textit{Método de mejora} que mejora la calidad de las soluciones.

\item Un \textit{Método de actualización del conjunto de referencia} que construye el conjunto de referencia con soluciones buenas y diversas.

\item Un \textit{Método de generación de subconjuntos} que determina cuales soluciones del conjunto de referencia sirven como base para la creación de nuevas soluciones.

\item Un \textit{Método de combinación de soluciones} que genera nuevas soluciones a partir de los subconjuntos creados por el método de generación de subconjuntos.

\end{enumerate}

El pseudocódigo general de SS se puede ver en el \textbf{Algoritmo~\ref{alg:SCA1}}

\begin{code}[includerangemarker=false,frame=single,label=alg:SCA1,caption=Pseudocódigo de Búsqueda en Dispersión,firstnumber=100, mathescape]
$Metodo$ $de$ $Generacion$ $y$ $Diversificacion$ - Crea un conjunto de soluciones diversas $P$ de tamano $Psize$

Construir el $RefSet$ con $b$ mejores y diversas soluciones de $P$, RefSet := $\lbrace S^1,...,S^b \rbrace$. Ordenarlas de manera creciente en base a su costo.  

$S_{mejor}$ := $S^1$
$NewSolutions$ := True
while ($NewSolutions$) do
	$NewSubset$ := GenerarSubconjuntos($RefSet$)
	$NewSolutions$ := False
    while ($NewSubset\neq\emptyset $) do
    	$Ss$ := Seleccionar($NewSubset$)
    	$S$ := Combinar(Ss) 
    	$x$ := Mejorar(S)
		if(($x$ $\notin$ $RefSet$) $\wedge$ ($Costo(x)<Costo(S^b$)))
			$S^b$ := $x$
			Ordenar($RefSet$)
			$NewSolutions$ := True											
		end if
		Eliminar($Ss$, $NewSubset$)    	    				
	end while
end while
$S_{mejor}$ := $S^1$

$\textbf{return}$ $S_{mejor}$
\end{code}


\subsection{Algoritmo Genético} \label{subsect:GA}

El \emph{algoritmo genético} (GA por sus siglas en inglés de \emph{Genetic Algorithm}) es una metaheurística poblacional utilizada por primera vez por John Holland en el año 1975. El enfoque de GA representa una analogía biológica en donde se simula el proceso evolutivo de una \textit{población} que refiere a un conjunto donde cada individuo representa una solución del problema. 

Cada solución perteneciente a la \textit{población} posee un valor asociado que recibe el nombre de \textit{fitness}. Éste determina si una solución es considerada como buena o no y aquellas que tengan mayor \textit{fitness} tendrán mayor oportunidad para aparearse y sobrevivir al paso del tiempo, lo cual permite que la población evolucione y obtenga cada vez mejores soluciones. Lo población inicial es generada de manera aleatoria, luego el algoritmo posee tres procedimientos fundamentales:

\begin{enumerate}
 
\item\textit{Selección:} Se encarga de elegir un subconjunto de la población, tal que las soluciones de mayor \textit{fitness} o calidad tengan mayor oportunidad de reproducirse. 

\item\textit{Cruce:} Se encarga de simular la reproducción de dos soluciones dentro de la población, las cuales para el efecto del GA serán llamadas ``padres''. El cruce da como resultado una o más soluciones que poseen características en común de ambos padres. Estas soluciones reciben el nombre de ``hijos'' para efectos del GA.

\item\textit{Mutación:} realiza una leve modificación de las características de una solución. Este proceso se utiliza para diversificar una nueva población sin necesidad de depender exclusivamente del proceso de cruce. 

\end{enumerate}

El pseudocódigo general de GA se puede ver en el \textbf{Algoritmo~\ref{alg:GA}}.
 
\begin{code}[includerangemarker=false,frame=single,label=alg:GA,caption=Pseudocódigo de Algoritmo Genético,firstnumber=100, mathescape]
t := 0
$P_0$ := GenerarPoblacionInicial()
Evaluar($P_t$)
while ($\neg$criterio_parada) do
	t := t+1
	$P_t$ := Seleccionar($P_{t-1}$)
	Cruzar($P_{t}$)
	Mutar($P_t$)
	Evaluar($P_t$)	
end while
$S$ := SeleccionarMejor($P_t$)

$\textbf{return}$ $S$
\end{code}
% Resultados
\chapter{Resultados} \label{chap:resultados}

A continuación se muestran los resultados de las pruebas realizadas.

\section{Resultados de las pruebas unitarias sobre el Módulo de Determinación de Origen de la Incompletitud.} 

\section{Resultados de las pruebas unitarias sobre el Módulo de Búsqueda Focalizada de Información}

A continuación se muestran los resultados de las pruebas unitarias realizadas sobre el Preprocesador de Textos y sobre el Extractor Focalizado de Información.

\subsection{Resultados de las pruebas unitarias del Preprocesador de Textos}

\begin{table}[h]
\caption{Resultados detallados la evaluación del Preprocesador de Textos}
\centering
\scriptsize
\begin{tabular*}{1\textwidth}{@{\extracolsep{\fill}} !{\vrule width 1pt} c !{\vrule width 1pt} c | c | c!{\vrule width 1pt} c | c | c!{\vrule width 1pt}}
\hline
Dominio & \multicolumn{3}{c!{\vrule width 1pt}}{\bf{Aprobados}} & \multicolumn{3}{c |}{\bf{No aprobados}}\\
\hline
 & Correctos & Con texto de más & \bf{Total aprobados} & Incompletos & Incorrectos & \bf{Total no aprobados}\\
\hline
Designaciones & 87.69\% & 7.69\% & \bf{95.39\%} & 3.07\% & 1.54\% & \bf{4.61\%}\\
\hline
Escalafón & 80.51\% & 12.98\% & \bf{93.6\%}  & 6.4\% & 0\% & \bf{6.4\%} \\
\hline
Jurados de Ascenso & 77.55\% & 16.32\% & \bf{93.88\%} & 0\% & 6.12\% & \bf{6.12\%} \\
\hline
\end{tabular*}
\label{tabla-resultados-preprocesamientoDatosDesignacion}

\end{table}

\subsection{Resultados de las pruebas unitarias del Extractor Focalizado}

En el apéndice \ref{chap:apendiceb} se muestra la totalidad de los resultados obtenidos de las pruebas unitarias realizadas sobre el Extractor Focalizado. Con el objetivo de seleccionar resultados ilustrativos, para estas pruebas se ha decidido mostrar para cada dominio de información los resultados obtenidos con un valor del Unit Hit Measure de 0.66, mostrando los resultados obtenidos con las probabilidades de campo faltante de 0; 0.33 y 0.66. El motivo de esto es que se decidió que dicho valor del Unit Hit Measure era representativo de los resultados obtenidos. Se consideró tomar el promedio de los 3 resultados obtenidos según las probabilidades de campo faltante pero se decidió que el efecto de dichas probabilidades era bastante relevante e importante, y por lo tanto conveniente de destacar.\\

De igual forma, para ilustrar el efecto del valor del Unit Hit Measure se decidió colocar los resultados obtenidos para valores de Unit Hit Measure de 0; 0.33 y 0.66 para el dominio de Designaciones.\\

TABLAS GO HERE\\

%DESIGNACIONES
%------------------------

\begin{table}
\centering
\caption{ Resultados de la evaluación del Extractor Focalizado - Dominio: Designaciones. UnitHit Measure mínimo:1.0}
\centering
\scriptsize
\begin{tabular*}{1\textwidth}{@{\extracolsep{\fill}} !{\vrule width 1pt} c !{\vrule width 1pt} c !{\vrule width 1pt} c | c | c !{\vrule width 1pt} c | c | c !{\vrule width 1pt}}
\hline
Campo & Prob. & \multicolumn{3}{c!{\vrule width 1pt}}{\bf{P. Aprobadas}} & \multicolumn{3}{c!{\vrule width 1pt}}{\bf{P. No aprobadas}}\\
\hline
\multicolumn{2}{!{\vrule width 1pt}c|}{ } & Correctos & Con texto de más & Aprobados & Incompletos & Incorrectos & No aprobados\\
\hline
\multirow{3}{*}{calificacion} 

	& 0.0
	& 74,86\% & 9,5\% & \bf{84,36\%} & 12,85\% & 2,79\% & \bf{15,64\%} \\
	\cline{3-8}

	& 0.33
	& 65,92\% & 7,26\% & \bf{73,18\%} & 24,58\% & 2,23\% & \bf{26,82\%} \\
	\cline{3-8}

	& 0.66
	& 39,11\% & 5,59\% & \bf{44,69\%} & 54,19\% & 1,12\% & \bf{55,31\%} \\
	\cline{3-8}

\hline
\multirow{3}{*}{fechaAsignacion} 

	& 0.0
	& 84,92\% & 1,12\% & \bf{86,03\%} & 13,41\% & 0,56\% & \bf{13,97\%} \\
	\cline{3-8}

	& 0.33
	& 73,18\% & 1,68\% & \bf{74,86\%} & 24,58\% & 0,56\% & \bf{25,14\%} \\
	\cline{3-8}

	& 0.66
	& 44,13\% & 1,12\% & \bf{45,25\%} & 54,75\% & 0\% & \bf{54,75\%} \\
	\cline{3-8}

\hline
\multirow{3}{*}{fechaFinal} 

	& 0.0
	& 96,09\% & 1,68\% & \bf{97,77\%} & 2,23\% & 0\% & \bf{2,23\%} \\
	\cline{3-8}

	& 0.33
	& 94,41\% & 3,91\% & \bf{98,32\%} & 1,68\% & 0\% & \bf{1,68\%} \\
	\cline{3-8}

	& 0.66
	& 87,15\% & 9,5\% & \bf{96,65\%} & 3,35\% & 0\% & \bf{3,35\%} \\
	\cline{3-8}

\hline
\multirow{3}{*}{motivo} 

	& 0.0
	& 88,83\% & 7,82\% & \bf{96,65\%} & 2,79\% & 0,56\% & \bf{3,35\%} \\
	\cline{3-8}

	& 0.33
	& 78,77\% & 13,97\% & \bf{92,74\%} & 6,7\% & 0,56\% & \bf{7,26\%} \\
	\cline{3-8}

	& 0.66
	& 63,69\% & 17,88\% & \bf{81,56\%} & 18,44\% & 0\% & \bf{18,44\%} \\
	\cline{3-8}

\hline
\multirow{3}{*}{Nombre} 

	& 0.0
	& 86,03\% & 2,79\% & \bf{88,83\%} & 10,61\% & 0,56\% & \bf{11,17\%} \\
	\cline{3-8}

	& 0.33
	& 70,95\% & 2,23\% & \bf{73,18\%} & 26,82\% & 0\% & \bf{26,82\%} \\
	\cline{3-8}

	& 0.66
	& 47,49\% & 2,79\% & \bf{50,28\%} & 49,16\% & 0,56\% & \bf{49,72\%} \\
	\cline{3-8}

\hline
\end{tabular*}
\label{Resultados-tabla-resultados-EFDesignaciones1.0}
\\

\end{table}


\begin{table}
\centering
\caption{ Resultados de la evaluación del Extractor Focalizado - Dominio: Designaciones. UnitHit Measure mínimo:0.66}
\centering
\scriptsize
\begin{tabular*}{1\textwidth}{@{\extracolsep{\fill}} !{\vrule width 1pt} c !{\vrule width 1pt} c !{\vrule width 1pt} c | c | c !{\vrule width 1pt} c | c | c !{\vrule width 1pt}}
\hline
Campo & Prob. & \multicolumn{3}{c!{\vrule width 1pt}}{\bf{P. Aprobadas}} & \multicolumn{3}{c!{\vrule width 1pt}}{\bf{P. No aprobadas}}\\
\hline
\multicolumn{2}{!{\vrule width 1pt}c|}{ } & Correctos & Con texto de más & Aprobados & Incompletos & Incorrectos & No aprobados\\
\hline
\multirow{3}{*}{calificacion} 

	& 0.0
	& 79,33\% & 9,5\% & \bf{88,83\%} & 8,38\% & 2,79\% & \bf{11,17\%} \\
	\cline{3-8}

	& 0.33
	& 64,8\% & 8,38\% & \bf{73,18\%} & 24,02\% & 2,79\% & \bf{26,82\%} \\
	\cline{3-8}

	& 0.66
	& 39,66\% & 4,47\% & \bf{44,13\%} & 54,75\% & 1,12\% & \bf{55,87\%} \\
	\cline{3-8}

\hline
\multirow{3}{*}{fechaAsignacion} 

	& 0.0
	& 83,8\% & 2,79\% & \bf{86,59\%} & 12,85\% & 0,56\% & \bf{13,41\%} \\
	\cline{3-8}

	& 0.33
	& 74,86\% & 2,79\% & \bf{77,65\%} & 22,35\% & 0\% & \bf{22,35\%} \\
	\cline{3-8}

	& 0.66
	& 50,28\% & 2,23\% & \bf{52,51\%} & 46,93\% & 0,56\% & \bf{47,49\%} \\
	\cline{3-8}

\hline
\multirow{3}{*}{fechaFinal} 

	& 0.0
	& 96,09\% & 2,79\% & \bf{98,88\%} & 1,12\% & 0\% & \bf{1,12\%} \\
	\cline{3-8}

	& 0.33
	& 92,18\% & 6,15\% & \bf{98,32\%} & 1,68\% & 0\% & \bf{1,68\%} \\
	\cline{3-8}

	& 0.66
	& 85,47\% & 10,61\% & \bf{96,09\%} & 3,91\% & 0\% & \bf{3,91\%} \\
	\cline{3-8}

\hline
\multirow{3}{*}{motivo} 

	& 0.0
	& 84,92\% & 12,85\% & \bf{97,77\%} & 1,68\% & 0,56\% & \bf{2,23\%} \\
	\cline{3-8}

	& 0.33
	& 72,63\% & 20,11\% & \bf{92,74\%} & 7,26\% & 0\% & \bf{7,26\%} \\
	\cline{3-8}

	& 0.66
	& 64,25\% & 20,11\% & \bf{84,36\%} & 15,64\% & 0\% & \bf{15,64\%} \\
	\cline{3-8}

\hline
\multirow{3}{*}{Nombre} 

	& 0.0
	& 89,39\% & 2,79\% & \bf{92,18\%} & 7,26\% & 0,56\% & \bf{7,82\%} \\
	\cline{3-8}

	& 0.33
	& 68,16\% & 2,79\% & \bf{70,95\%} & 28,49\% & 0,56\% & \bf{29,05\%} \\
	\cline{3-8}

	& 0.66
	& 40,78\% & 2,23\% & \bf{43,02\%} & 56,98\% & 0\% & \bf{56,98\%} \\
	\cline{3-8}

\hline
\end{tabular*}
\label{Resultados-tabla-resultados-EFDesignaciones0.66}
\\

\end{table}


\begin{table}
\centering
\caption{ Resultados de la evaluación del Extractor Focalizado - Dominio: Designaciones. UnitHit Measure mínimo:0.33}
\centering
\scriptsize
\begin{tabular*}{1\textwidth}{@{\extracolsep{\fill}} !{\vrule width 1pt} c !{\vrule width 1pt} c !{\vrule width 1pt} c | c | c !{\vrule width 1pt} c | c | c !{\vrule width 1pt}}
\hline
Campo & Prob. & \multicolumn{3}{c!{\vrule width 1pt}}{\bf{P. Aprobadas}} & \multicolumn{3}{c!{\vrule width 1pt}}{\bf{P. No aprobadas}}\\
\hline
\multicolumn{2}{!{\vrule width 1pt}c|}{ } & Correctos & Con texto de más & Aprobados & Incompletos & Incorrectos & No aprobados\\
\hline
\multirow{3}{*}{calificacion} 

	& 0.0
	& 80,45\% & 9,5\% & \bf{89,94\%} & 7,26\% & 2,79\% & \bf{10,06\%} \\
	\cline{3-8}

	& 0.33
	& 63,69\% & 8,94\% & \bf{72,63\%} & 25,14\% & 2,23\% & \bf{27,37\%} \\
	\cline{3-8}

	& 0.66
	& 40,78\% & 5,59\% & \bf{46,37\%} & 53,07\% & 0,56\% & \bf{53,63\%} \\
	\cline{3-8}

\hline
\multirow{3}{*}{fechaAsignacion} 

	& 0.0
	& 83,24\% & 3,35\% & \bf{86,59\%} & 12,85\% & 0,56\% & \bf{13,41\%} \\
	\cline{3-8}

	& 0.33
	& 74,3\% & 2,79\% & \bf{77,09\%} & 22,35\% & 0,56\% & \bf{22,91\%} \\
	\cline{3-8}

	& 0.66
	& 49,72\% & 3,35\% & \bf{53,07\%} & 46,93\% & 0\% & \bf{46,93\%} \\
	\cline{3-8}

\hline
\multirow{3}{*}{fechaFinal} 

	& 0.0
	& 82,68\% & 16,2\% & \bf{98,88\%} & 1,12\% & 0\% & \bf{1,12\%} \\
	\cline{3-8}

	& 0.33
	& 58,1\% & 40,78\% & \bf{98,88\%} & 1,12\% & 0\% & \bf{1,12\%} \\
	\cline{3-8}

	& 0.66
	& 65,36\% & 30,73\% & \bf{96,09\%} & 3,91\% & 0\% & \bf{3,91\%} \\
	\cline{3-8}

\hline
\multirow{3}{*}{motivo} 

	& 0.0
	& 59,78\% & 37,99\% & \bf{97,77\%} & 1,68\% & 0,56\% & \bf{2,23\%} \\
	\cline{3-8}

	& 0.33
	& 49,16\% & 44,69\% & \bf{93,85\%} & 5,59\% & 0,56\% & \bf{6,15\%} \\
	\cline{3-8}

	& 0.66
	& 53,07\% & 31,28\% & \bf{84,36\%} & 15,08\% & 0,56\% & \bf{15,64\%} \\
	\cline{3-8}

\hline
\multirow{3}{*}{Nombre} 

	& 0.0
	& 90,5\% & 2,79\% & \bf{93,3\%} & 6,15\% & 0,56\% & \bf{6,7\%} \\
	\cline{3-8}

	& 0.33
	& 72,07\% & 1,68\% & \bf{73,74\%} & 25,7\% & 0,56\% & \bf{26,26\%} \\
	\cline{3-8}

	& 0.66
	& 39,11\% & 1,68\% & \bf{40,78\%} & 58,66\% & 0,56\% & \bf{59,22\%} \\
	\cline{3-8}

\hline
\end{tabular*}
\label{Resultados-tabla-resultados-EFDesignaciones0.33}
\\

\end{table}


\begin{table}
\centering
\caption{ Resultados de la evaluación del Extractor Focalizado - Dominio: Escalafon. UnitHit Measure mínimo:0.66}
\centering
\scriptsize
\begin{tabular*}{1\textwidth}{@{\extracolsep{\fill}} !{\vrule width 1pt} c !{\vrule width 1pt} c !{\vrule width 1pt} c | c | c !{\vrule width 1pt} c | c | c !{\vrule width 1pt}}
\hline
Campo & Prob. & \multicolumn{3}{c!{\vrule width 1pt}}{\bf{P. Aprobadas}} & \multicolumn{3}{c!{\vrule width 1pt}}{\bf{P. No aprobadas}}\\
\hline
\multicolumn{2}{!{\vrule width 1pt}c|}{ } & Correctos & Con texto de más & Aprobados & Incompletos & Incorrectos & No aprobados\\
\hline
\multirow{3}{*}{EsAscendido.Nombre} 

	& 0.0
	& 92,31\% & 1,54\% & \bf{93,85\%} & 5,38\% & 0,77\% & \bf{6,15\%} \\
	\cline{3-8}

	& 0.33
	& 76,92\% & 1,54\% & \bf{78,46\%} & 20,77\% & 0,77\% & \bf{21,54\%} \\
	\cline{3-8}

	& 0.66
	& 43,08\% & 0\% & \bf{43,08\%} & 56,15\% & 0,77\% & \bf{56,92\%} \\
	\cline{3-8}

\hline
\multirow{3}{*}{EsAscendido.NombreTrabajo} 

	& 0.0
	& 94,62\% & 1,54\% & \bf{96,15\%} & 2,31\% & 1,54\% & \bf{3,85\%} \\
	\cline{3-8}

	& 0.33
	& 70,77\% & 1,54\% & \bf{72,31\%} & 26,92\% & 0,77\% & \bf{27,69\%} \\
	\cline{3-8}

	& 0.66
	& 43,08\% & 1,54\% & \bf{44,62\%} & 53,85\% & 1,54\% & \bf{55,38\%} \\
	\cline{3-8}

\hline
\multirow{3}{*}{EsAscendido.escalafon} 

	& 0.0
	& 97,69\% & 0\% & \bf{97,69\%} & 2,31\% & 0\% & \bf{2,31\%} \\
	\cline{3-8}

	& 0.33
	& 90,77\% & 0\% & \bf{90,77\%} & 9,23\% & 0\% & \bf{9,23\%} \\
	\cline{3-8}

	& 0.66
	& 71,54\% & 0\% & \bf{71,54\%} & 28,46\% & 0\% & \bf{28,46\%} \\
	\cline{3-8}

\hline
\multirow{3}{*}{EsAscendido.fecha} 

	& 0.0
	& 95,38\% & 2,31\% & \bf{97,69\%} & 2,31\% & 0\% & \bf{2,31\%} \\
	\cline{3-8}

	& 0.33
	& 76,15\% & 9,23\% & \bf{85,38\%} & 14,62\% & 0\% & \bf{14,62\%} \\
	\cline{3-8}

	& 0.66
	& 60,77\% & 8,46\% & \bf{69,23\%} & 30,77\% & 0\% & \bf{30,77\%} \\
	\cline{3-8}

\hline
\multirow{3}{*}{EsAscendido.postergado} 

	& 0.0
	& 99,23\% & 0\% & \bf{99,23\%} & 0,77\% & 0\% & \bf{0,77\%} \\
	\cline{3-8}

	& 0.33
	& 96,15\% & 2,31\% & \bf{98,46\%} & 1,54\% & 0\% & \bf{1,54\%} \\
	\cline{3-8}

	& 0.66
	& 84,62\% & 12,31\% & \bf{96,92\%} & 3,08\% & 0\% & \bf{3,08\%} \\
	\cline{3-8}

\hline
\end{tabular*}
\label{Resultados-tabla-resultados-EFEscalafon0.66}
\\

\end{table}

\begin{table}
\centering
\caption{ Resultados de la evaluación del Extractor Focalizado - Dominio: JuradosAscenso. UnitHit Measure mínimo:0.66}
\centering
\scriptsize
\begin{tabular*}{1\textwidth}{@{\extracolsep{\fill}} !{\vrule width 1pt} c !{\vrule width 1pt} c !{\vrule width 1pt} c | c | c !{\vrule width 1pt} c | c | c !{\vrule width 1pt}}
\hline
Campo & Prob. & \multicolumn{3}{c!{\vrule width 1pt}}{\bf{P. Aprobadas}} & \multicolumn{3}{c!{\vrule width 1pt}}{\bf{P. No aprobadas}}\\
\hline
\multicolumn{2}{!{\vrule width 1pt}c|}{ } & Correctos & Con texto de más & Aprobados & Incompletos & Incorrectos & No aprobados\\
\hline
\multirow{3}{*}{Jurado.MiembroPrincipalExterno} 

	& 0.0
	& 90,37\% & 0\% & \bf{90,37\%} & 2,22\% & 7,41\% & \bf{9,63\%} \\
	\cline{3-8}

	& 0.33
	& 88,89\% & 0,74\% & \bf{89,63\%} & 2,96\% & 7,41\% & \bf{10,37\%} \\
	\cline{3-8}

	& 0.66
	& 75,56\% & 7,41\% & \bf{82,96\%} & 12,59\% & 4,44\% & \bf{17,04\%} \\
	\cline{3-8}

\hline
\multirow{3}{*}{Jurado.MiembroPrincipalInterno} 

	& 0.0
	& 97,78\% & 0\% & \bf{97,78\%} & 2,22\% & 0\% & \bf{2,22\%} \\
	\cline{3-8}

	& 0.33
	& 97,04\% & 0\% & \bf{97,04\%} & 2,96\% & 0\% & \bf{2,96\%} \\
	\cline{3-8}

	& 0.66
	& 81,48\% & 1,48\% & \bf{82,96\%} & 17,04\% & 0\% & \bf{17,04\%} \\
	\cline{3-8}

\hline
\multirow{3}{*}{Jurado.Presidente} 

	& 0.0
	& 97,78\% & 0\% & \bf{97,78\%} & 2,22\% & 0\% & \bf{2,22\%} \\
	\cline{3-8}

	& 0.33
	& 97,78\% & 0\% & \bf{97,78\%} & 2,22\% & 0\% & \bf{2,22\%} \\
	\cline{3-8}

	& 0.66
	& 82,96\% & 0\% & \bf{82,96\%} & 17,04\% & 0\% & \bf{17,04\%} \\
	\cline{3-8}

\hline
\multirow{3}{*}{Jurado.SuplenteExterno} 

	& 0.0
	& 94,81\% & 0\% & \bf{94,81\%} & 2,22\% & 2,96\% & \bf{5,19\%} \\
	\cline{3-8}

	& 0.33
	& 91,85\% & 2,96\% & \bf{94,81\%} & 2,22\% & 2,96\% & \bf{5,19\%} \\
	\cline{3-8}

	& 0.66
	& 77,04\% & 6,67\% & \bf{83,7\%} & 14,07\% & 2,22\% & \bf{16,3\%} \\
	\cline{3-8}

\hline
\multirow{3}{*}{Jurado.SuplenteInterno} 

	& 0.0
	& 95,56\% & 2,22\% & \bf{97,78\%} & 2,22\% & 0\% & \bf{2,22\%} \\
	\cline{3-8}

	& 0.33
	& 95,56\% & 2,22\% & \bf{97,78\%} & 2,22\% & 0\% & \bf{2,22\%} \\
	\cline{3-8}

	& 0.66
	& 74,81\% & 2,96\% & \bf{77,78\%} & 22,22\% & 0\% & \bf{22,22\%} \\
	\cline{3-8}

\hline
\multirow{3}{*}{Profesor.Departamento} 

	& 0.0
	& 97,78\% & 0\% & \bf{97,78\%} & 2,22\% & 0\% & \bf{2,22\%} \\
	\cline{3-8}

	& 0.33
	& 97,04\% & 0\% & \bf{97,04\%} & 2,96\% & 0\% & \bf{2,96\%} \\
	\cline{3-8}

	& 0.66
	& 83,7\% & 0\% & \bf{83,7\%} & 16,3\% & 0\% & \bf{16,3\%} \\
	\cline{3-8}

\hline
\multirow{3}{*}{Profesor.Nombre} 

	& 0.0
	& 97,78\% & 0\% & \bf{97,78\%} & 1,48\% & 0,74\% & \bf{2,22\%} \\
	\cline{3-8}

	& 0.33
	& 96,3\% & 0\% & \bf{96,3\%} & 2,96\% & 0,74\% & \bf{3,7\%} \\
	\cline{3-8}

	& 0.66
	& 77,78\% & 0\% & \bf{77,78\%} & 21,48\% & 0,74\% & \bf{22,22\%} \\
	\cline{3-8}

\hline
\multirow{3}{*}{Trabajo.Escalafon} 

	& 0.0
	& 98,52\% & 0\% & \bf{98,52\%} & 1,48\% & 0\% & \bf{1,48\%} \\
	\cline{3-8}

	& 0.33
	& 97,04\% & 0\% & \bf{97,04\%} & 2,96\% & 0\% & \bf{2,96\%} \\
	\cline{3-8}

	& 0.66
	& 87,41\% & 0\% & \bf{87,41\%} & 12,59\% & 0\% & \bf{12,59\%} \\
	\cline{3-8}

\hline
\multirow{3}{*}{Trabajo.Nombre} 

	& 0.0
	& 97,78\% & 0\% & \bf{97,78\%} & 1,48\% & 0,74\% & \bf{2,22\%} \\
	\cline{3-8}

	& 0.33
	& 97,78\% & 0\% & \bf{97,78\%} & 1,48\% & 0,74\% & \bf{2,22\%} \\
	\cline{3-8}

	& 0.66
	& 77,78\% & 0\% & \bf{77,78\%} & 21,48\% & 0,74\% & \bf{22,22\%} \\
	\cline{3-8}

\hline
\end{tabular*}
\label{Resultados-tabla-resultados-EFJuradosAscenso0.66}
\\

\end{table}

\subsection{Análisis de los resultados de las pruebas unitarias del Buscador Focalizado de Información}

Los resultados obtenidos varían muchísimo según las tres variables que entran en juego en este experimento: el dominio de información (al cual está asociado una semántica que determina a su vez un extractor especificado por expresiones regulares), el Unit Hit Measure (UHM) y la probabilidad (P) de que cada uno de los campos falten.\\

Puede verse en primera instancia como algunos campos obtenidos tuvieron una aprobación considerablemente alta: postergado en el Dominio Escalafón obtuvo un 98.46\%, escalafón en el caso de los jurados de ascenso 97,04\% y fecha de Designación en el caso del dominio Designación 98,88\%. De igual forma, hay campos con precisión muchísimos menor: nombre del Miembro Principal Externo en el caso de Jurados de Ascenso obtuvó 89,63\%, el nombre del profesor ascendido en el caso de Escalafón 78,46\%, la calificación de la designación (el nombre del cargo y la unidad de la designación) 73,18\%. \\

Todos estos valores son representativos de cada uno de los 3 dominios y se corresponden a las pruebas hechas con un UHM de 0.66 y P de 0,33. Esto parámetros son considerados representativos y válidos para sacar algunas conclusiones generales. Los resultados obtenidos reflejan muchos de los retos que supone trabajar con expresiones regulares, así como lo útil que es en el ámbito de FIE el contexto de extracción.  En lo sucesivo se explicará brevemente los resultados obtenidos y el comportamiento del Extractor Focalizado. Al menos de que se indique lo contrario, al hablar de resultados puntuales se hará referencia a los resultados obtenidos con UHM de 0,66 y P de ,33. \\

El primer elemento que debe considerarse para entender los resultados es el funcionamiento de las expresiones regulares. Las expresiones regulares utilizadas funcionan especificando delimitadores anteriores y posteriores al valor que se quiere extraer. Dichos delimitadores son muchos casos cadenas de texto estáticas elegidas por un usuario experto luego de analizar el documento y encontrar repeteciones de frases. En muchos casos se especifican en lugar de cadenas de texto clases de caracteres que especifican un conjunto de símbolos que puede repetirse. La mayoría de las veces se busca combinar ambas posibilidades y que los delimitadores estén basados en las cadenas de texto modeladas con clases de caracteres para minimizar los errores de tipeo, de mayúsculas y minúsculas, diferencias en el formato, etcétera. \\

En todo caso lo relevante en este punto, más allá de entender el proceso que tiene que seguir un usuario experto para configurar el extractor, es entender que a pesar de la riqueza que dan las clases de caracteres,  las expresiones regulares pueden considerarse como reglas bastante rígidas y estáticas. Como no son el producto de un proceso de aprendizaje de máquina o de un mecanismo probabilístico, se basan en la configuración que haga el usuario y dependen muchísimo de que la regularidad observada al hacer la configuración sea generalizable a todos los documentos. \\

Esto es complicado en la medida en la que se haga más amplio el conjunto de documentos sobre los cuales se hace la extracción: a menudo la nomenclatura utilizada, el estilo de redacción, el formato de los documentos, el posicionamiento de la información dentro del documento cambia de un año a otro o inclusive de un mes a otro. La labor del usuario experto es generalizar al máximo las reglas con las que configura el extractor, pero en muchos casos es una tarea cercana a lo imposible.\\

En muchos casos es imposible generalizar los delimitadores y se depende del conocimiento del usuario experto y de los documentos utilizados para modelar el documento. Por ejemplo, el campo nombre de la designación es uno de los más dificiles para generalizar: el resultado obtenido es uno de los más bajos (70,95\%). El motivo de esto es que los delimitadores utilizados para definir este campo cambian muchísimo y es poco práctico compilar una lista con todos los posibles nombrse que puede tener una persona. \\

Por otro lado, hay campos que dada la regularidad de sus valores es fácil de modelar. Es el caso del nombre del escalafón en el dominio de Jurados de trabajo de Ascenso (que tiene 3 valores posibles: titular, asociado y agregado), de las fechas (que se presentan típicamente en formatos dd/MM/aaaa, dd-mm-aaaa), si el ascenso es postergado o no (típicamente se identificamente con una frase del tipo: veredicto desfavorable, veredicto reprobatorio, rechazar dicho trabajo, entre otras). Dichos campos obtuvieron precisiones de 97.04\%, 87.11 (promedio de los 3 campos de tipo fecha presentes en los 3 dominios), 98.46\% respectivamente. Nótese sin embargo, que la precisión de estos campos no es perfecta. En muchos casos una examinación de los resultados puntuales en estos campos indican errores de tipeo que explican los errores del orden del 1-2\%. En otros casos, como en las fechas, existen documentos en los que la fecha se escribe extensivamente en lugar de la sintaxis abreviada (15 de diciembre de mil novecientos noventa y nueve, por ejemplo, en lugar de 15-12-1999). \\

Otra observación que vale la pena hacer es el efecto que puede tener la estructura del documento y el posicionamiento de los campos dentro del mismo. Esto se puede apreciar en el dominio de designación de jurados de ascenso: los nombres de los miembros del jurado, campos del mismo tipo y que tienen todos el mismo formato (en esencia, el nombre), tienen resultados diferentes según el posicionamiento que tienen dentro del texto. El posicionamiento influye precisamente porque determina la selección de los delimitadores del documento: hay veces que el usuario que construye las expresiones regulares puede beneficiarse de que un campo esté de último en el párrafo o que esté justo después de una palabra clave. \\

En el caso concreto de los jurados, el resultado promedio de los nombres de los 5 jurados es de 94.8. El resultado del Miembro Principal Exterbi es sin embargo 89,63. Este resultado puede atribuirse a que las expresiones regulares óptimas para este campo fueron construidas esperando que el Miembro Principal Externo sea el último miembro a definirse; en los casos en los que no hay un miembro principal externo el último miembro es el miembro principal interno y por lo tanto se tiene una respuesta equivocada. \\

Hay también varios casos en los que delimitadores de una unidad de documento tienen errores. Como bien se ha dicho, una unidad de documento es un fragmento de texto que se refiere a una designación, ascenso o jurado de ascenso. Si se tiene más de una unidad en lo que el extractor reconoce como una (por problemas en los delimitadores de las mismas), se pierde toda la información de las unidades adicionales. Esto demuestra lo crítico que es construir las mejores expresiones regulares para delimitar unidades de documento. Quizás sea interesante probar otro mecanismo que no dependa de la existencia de unidades (look ahead, entre otros).  \\

En otro casos, el extractor no permite procesar corrrecciones registradas en los documentos. Es decir, puede ocurrir que en un acta se especifiquen correcciones sobre la información de una designación, ascenso o jurado que fue anunciado en un acta previa. De cierta manera, se tienen más de una unidad de documento que se refieren a un mismo hecho: la designación, el ascenso o el jurado inicial y las correcciones a los mismos. El extractor es incapaz de diferenciar entre ambos, aunque ambos tengan un UHM alto, y se queda con el primero que consiga. Esto podría solucionarse mejorando el extractor, de forma que pueda identificar cuando una unidad está vinculada a otra unidad presente en el conjunto de los documentos. Esta tarea no puede resolverse exclusivamente escribiendo mejores expresiones regulares: se tiene que introducir una forma de discriminar entre unidades con igual Unit Hit Measure.\\

De igual forma otro problema asociado al procesamiento por unidades de documento fue encontrado en el dominio de Designaciones. Se encontraron muchísimos casos en los que hay multiples designaciones para una misma persona para un mismo día, y que por lo tanto es imposible diferenciar como unidades diferentes. Esto es, si se tiene un dominio cuya naturaleza hace que los contextos de extracción sean iguales entre unidades de documentos que se refieren a cosas diferentes, se pierde información. Esta hipótesis sirve para explicar los resultados obtenidos en la calificación de la designación: 73,18\%. \\

Los resultados obtenidos son de esta manera bastante interesantes y la cantidad de información da para sacar muchísimas conclusiones sobre la complejidad que supone la tarea hacer ingeniería de extracción. La tecnología utilizada para hacer la extracción, expresiones regulares, supone retos adicionales y los resultados varían muchísimo según el campo. Es importante recordar que estos resultados son obtenidos en el contexto de documentos semi-estructurados y con una cierta regularidad en su estructura y redacción; es virtualmente imposible pretender hacer extracción en dominios con estilo de redacción más libre salvo para datos muy, muy puntuales. Esto es sin duda alguna uno de los motivos por los cuales las expresiones regulares han sido relegadas en el estudio de las áreas de Procesamiento de Lenguaje Natural y Extracción de Información. Hoy por hoy se apuesta más por métodos estocásticos y aprendizaje de máquina para realizar estas tareas. \\

Hasta ahora se ha analizado exclusivamente el comportamiento de las expresiones regulares según los dominios y los campos tomando en cuenta si el resultado obtenido es aprobado o no. Los resultados sin embargo pueden ser refinados en dos subcategorías dentro de las categorías aprobado y no aprobado. Dentro de la categoría aprobado se tiene que el resultado sea correcto (que sea un match exacto) o que la respuesta tenga texto de más (la respuesta obtenida contiene a la respuesta correcta). Dentro de la categoría no aprobado se tienen incompletos (la respuesta correcta contiene a la respuesta obtenido) e incompleto (las respuestas obtenidas y correctas son completamente disjuntos). \\

El motivo por el cual se decidió considerar estas categorías es que se considera ilustrativo e importante para profundizar el entendimiento de los resultados. Por ejemplo, el campo motivo de Designaciones registró un 20,11\% de respuestas aprobadas que, sin embargo, contienen texto de más. En este caso particular el motivo es que si bien es relativamente fácilmente conseguir el delimitador inicial del motivo de la designación, no es tan fácil conseguir el delimitador final. Por ende sucede a menudo que el extractor consigue el motivo pero captura texto adicional. \\

De igual forma hay varios campos que registraron porcentajes de incompletitud altos. El nombre del profesor en Designaciones registró un 28,49\% de incompletitud, lo cual quiere decir que si bien se extrajó parte del nombre, no se tuvo completo. El motivo de esto, nuevamente, es que dada la naturaleza del campo y de la posición del mismo en el documento, se hizó muy complicado elegir un delimitador final para el nombre que siempre lo capturara por completo. No es el caso de otros nombres de personas en otros dominios, que sí registraron tasas de acierto más deseables (nombres de los jurados, nombre del profesor que es ascendido, entre otros)\\

Por último, y no menos importante, la tasa de incorrectitud se refiere a los casos en los que la respuesta obtenida no tiene absolutamente nada que ver con la respuesta correcta (no contiene ni es contenida una por la otra). En casí todos los casos la proporción de respuestas es despreciable: en la mayoría de los casos es 0 (11 de 19),  el promedio de la proporción de incorrectitud de todos los campos 0,87\% y sólo en un campo pasa de 3\% (el nombre del miembro principal externo, con un 7,41\%). La existencia de estas respuestas incorrectas puede atribuirse a problemas en la selección de los delimitadores. \\

Queda en este punto comentar sobre el efecto que tiene la probabilidad de que falte un campo y el Hit Measure en los resultados obtenidos. Cómo se ha dicho anteriormente la probabilidad de que falte un campo es un elemento introducido en los experimentos para modelar el efecto que tiene la riqueza que tiene el contexto de extracción a la hora de hacer una búsqueda focalizada. Los archivos que se utilizaron para generar las pruebas contienen toda la información que se puede extraer de los documentos; sin embargo a la hora de hacer las pruebas, al incorporar el valor de cada campo se genera un número aleatorio; si es mayor que la probabilidad de la prueba se incorpora al contexto, de lo contrario no se incorpora. Esto permite simular condiciones normales de una extracción focalizada en la que no se tienen todos los valores de los campos. \\

El efecto de la probabilidad en los resultados obtenidos es bastante claro: se puede observar que conforme aumenta la probabilidad de que falte un campo (esto es, que el contexto de extracción sea menos rico) disminuye la proporción de resultados aprobados. Este comportamiento se puede observar en practicamente todos los campos y dominios. La disminución de la tasa de aprobación varía sin embargo bastante de acuerdo al campo: en algunos campos es del orden de las décimas mientras que en algunos otros es del orden de las decenas. Esto puede deberse a que existen campos que son más relevantes para poder discriminar entre unidades de documento. Si toca el caso de que el campo que se esté buscando es importante para identificar cada unidad, el efecto de que el contexto esté incompleto puede ser mayor. \\

Por último el UHM es una medida de qué tan exigente es el extractor focalizado al seleccionar unidades de documento. Esto es, mide que porcentaje de campos contenidos en el contexto de extracción están presentes en cada unidad de documento: 1.0 indica que todos los valores del contexto están en la unidad, .66 y .33 que el 66\% y el 33\% están. En general en la medida en la que se flexibiliza la exigencia del extractor con el UHM, aumenta el número de resultados aprobados. El motivo de esto es que relajar la exigencia sólo permite que en caso de no tener mejores unidades se procesen unidades con una menor probabilidad de que se refieran al contexto de extracción. El extractor, como se ha dicho, ordena las unidades encontradas de acuerdo al unit hit measure y busca extraer primero el valor del campo faltante de las unidades con mayor UHM. \\

La importante del Unit Hit Measure es que de no tener exigencias sobre el mínimo valor de cada unidad se pueden tener resultados completamente incorrectos si el extractor no ubica las unidades de documento correctas y pasa a analizar unidades con menor UHM. \\

A manera de conclusión, el análisis aquí busca ilustrar, más que desmenuzar toda la información que se puede extraer de estas pruebas, los principales hallazgos producto de la utilización del extractor. Los hallazgos más importantes son que la utilización de expresiones regulares depende, como podría esperarse, fuertemente de la regularidad de un documento en su estructura y estilo de redacción, que está condicionada a errores humanos de tipeo y que depende de los delimitadores que un usuario experto pueda definir. De igual forma, depende fuertemente de cada uno de los campos y de la forma como se encuentran esos campos en el texto. Influye también la posición relativa que tengan cada uno de ellos.\\

De igual manera, la existencia de información previa (el contexto de extracción, que es precisamente la principal ventaja de la extracción focalizada) incide determinantemente en el hallazgo de información faltante en la ontología.\\

% Conclusiones
\chapter*{Conclusiones y recomendaciones} \label{chap:conclusiones}
\addcontentsline{toc}{chapter}{Conclusiones y recomendaciones}

En este capítulo se presentan los hallazgos y contribuciones de este trabajo y se dan algunas recomendaciones.\\

Seis metaheurísticas del estado del arte fueron estudiadas y posteriormente implementadas para la resolución del Problema de Enrutamiento de Vehículos con Entrega y Recepción Simultánea. 
Se probaron para los dos tipos de instancias Dethloff y SalhiNagy, y se realizó un estudio comparativo con las metaheurísticas afines.

Para instancias Dethloff la metaheurística implementada GTS-M dió resultados aproximados a los mejores re\-por\-ta\-dos hasta el momento, reduciendo los tiempos de ejecución a un poco menos de la mitad. La metaheurística ILS-VND-M no logró superar los resultados publicados en el trabajo referenciado. La metaheurística GTS-M es altamente recomendada por su calidad de solución y menor tiempo.

Para instancias SalhiNagy las implementaciones de metaheurísticas de trayectoria (ILS-VND-M y GTS-M)  no lograron superar los resultados publicados en los trabajos afines. Esto puede deberse a una falta de mejor entonación de los parámetros de las metaheurísticas. La implementación de SS-M tampoco logró superar los resultados de su afín. La metaheurística implementada que generó mejores resultados fue AS-M, sin embargo consume casi cuatro veces el tiempo de ejecución de la metaheurística con mejores resultados (GTS). De las metaheurísticas implementadas AS-M es recomendada por sus resultados a\-proxi\-ma\-dos de calidad y tiempo.

Independientemente del tipo de instancia utilizada, las metaheurísticas de trayectoria presentan mejor comportamiento que las poblacionales.\\

Cada una de las metaheurísticas implementadas se hibridó con VND. El aporte de VND fue fundamental para lograr un mejoramiento de las metaheurísticas implementadas con respecto a los trabajos referenciados. La hibridación sirvió para mejorar los mejores costos promedio en la mayoría de los casos.\\

La entonación de los parámetros de las metaheurísticas es un proceso de vital importancia ya que de ella depende en gran medida la calidad de las soluciones encontradas. Los muestreos realizados fueron de ayuda para encontrar aquellos parámetros que influyan más en la calidad de las soluciones, así como para acotar los rangos de cada uno de los parámetros. Más adelante, las pruebas de entonación de parámetros lograron conseguir los valores de los parámetros que generaban, en promedio, las mejores soluciones.\\

En trabajos futuros se recomienda realizar de forma más exhaustiva la entonación de parámetros, realizando análisis de varianza a todos los parámetros, ya que el muestreo de los parámetros da como resultado sólo un estimador de la influencia de estos parámetros en el comportamiento de la metaheurística y no es una forma certera de conseguir los parámetros más influyentes, así como los rangos de entonación a utilizar. Adicionalmente, se recomienda la ejecución de las metaheurísticas utilizando otras clases de instancia, para de esta manera, asegurar el buen comportamiento de las metaheurísticas en la resolución de VRPSPD en casos de la vida real.

%\begin{table}[ht]
\caption{Resultados de la ejecución de la metaheurística GTS, utilizando instancias de Dethloff con la configuración XXX}
\centering
\begin{tabular}{c c c c c}
\hline\hline
Instancia & Costo mínimo & Tiempo(seg.) & Costo promedio & Tiempo promedio(seg.) \\ [0.5ex]
\hline
c0530.txt & 640.549071 & 1.19 & 640.549 & 1.196 \\
c0531.txt & 697.835518 & 1.06 & 698.368 & 1.56 \\
c0532.txt & 659.335981 & 1.74 & 659.336 & 1.572 \\
c0533.txt & 680.604414 & 1.09 & 682.989 & 1.374 \\
c0534.txt & 690.499806 & 1.75 & 693.912 & 2.472 \\
[1ex]\hline
\end{tabular}
\label{table:nonlin}
\end{table}
\begin{table}[ht]
\caption{Resultados de la ejecución de la metaheurística GTS, utilizando instancias de Dethloff con la configuración XXX}
\centering
\begin{tabular}{c c c c c}
\hline\hline
Instancia & Costo mínimo & Tiempo(seg.) & Costo promedio & Tiempo promedio(seg.) \\ [0.5ex]
\hline
c0530.txt & 636.336875 & 3.68 & 639.707 & 2.166 \\
c0531.txt & 697.835518 & 1.35 & 699.381 & 1.756 \\
c0532.txt & 659.335981 & 1.24 & 659.336 & 1.402 \\
c0533.txt & 680.041387 & 4.57 & 683.589 & 2.258 \\
c0534.txt & 690.499806 & 2.02 & 697.517 & 2.062 \\
[1ex]\hline
\end{tabular}
\label{table:nonlin}
\end{table}

\begin{table}[ht]
\caption{Resultados de la ejecución de la metaheurística GTS, utilizando instancias de SalhiNagy con la configuración XXX}
\centering
\begin{tabular}{c c c c c}
\hline\hline
Instancia & Costo mínimo & Tiempo(seg.) & Costo promedio & Tiempo promedio(seg.) \\ [0.5ex]
\hline
vrpnc1x.txt & 470.479372 & 3.06 & 475.716 & 2.376 \\
vrpnc1y.txt & 470.479372 & 2.34 & 473.529 & 2.536 \\
vrpnc2x.txt & 686.772205 & 5.43 & 691.376 & 5.348 \\
vrpnc2y.txt & 682.387217 & 3.03 & 690.264 & 4.622 \\
vrpnc3x.txt & 721.270834 & 14.85 & 725.903 & 11.458 \\
[1ex]\hline
\end{tabular}
\label{table:nonlin}
\end{table}

% Crea el glosario 
\printglossaries

% Establece las citas y bibliografia
\bibliographystyle{alpha.bst}
\bibliography{myrefs}

% Crea el apendice
\appendix

\end{document}
